\section{Multi‑axial atomicity as separatedness}\label{sec:separated}

\begin{definition}[Separated object]
An object $s \in \mathcal{A}$ is called \emph{separated} if for every object $X \not\cong s$, the hom‑set $\operatorname{Hom}_{\mathcal{A}}(s,X)$ is empty.
\end{definition}

\begin{theorem}[Separated object characterisation]\label{th:separated}
Let $o \in \mathcal{A}$.  $o$ is multi‑axially atomic if and only if $o$ is a separated object.
\end{theorem}
\begin{proof}
Immediate from Lemma~\ref{lem:atomic_separated}: multi‑axial atomicity is equivalent to the absence of decomposition edges out of $o$, hence to the absence of any non‑trivial path to another object.
\end{proof}

In categorical terms, multi‑axial atomicity is exactly the property of being \emph{indecomposable} with respect to the decomposition preorder: no further refinement is possible.
Such separated objects are precisely the maximal elements in the preorder induced by the free category $\mathcal{A}$.